% "Tips": shared by the rulebook and the platform manuals (docs/manual/*).
% Each document sets its own heading above this.
None of this is a rule. It is what experienced \Vp{player} tend to know.

\subsection{Early: get \Vp{hive} on the \V{grid}}
\begin{itemize}
  \item \textbf{Found \Vp{hive} whenever you can.} The \V{finderscut} is a free \V{coin}, and you can buy \Rule{maxbuy} more on the same
    turn, while the price is at its lowest. That is a head start on \V{root} that nobody else gets.
  \item \textbf{Small \Vp{hive} pay too.} A \V{hive} that gets \V{hivejack}ed early still pays its \V{root} and \V{sudo}.
    Don't ignore a \V{hive} just because it will never be big.
  \item \textbf{Don't stretch yourself thin.} With only \Rule{maxbuy} \Vp{coin} a turn, being \V{root} of one syndicate is usually worth
    more than being third in three.
\end{itemize}

\begin{example}[: A QUICK FLIP]
You found a low-tier \V{hive} of 2 \Vp{node}, take your \V{finderscut}, and buy 3 more \V{coin} at \Y{200} each (\Y{600}).
Nobody else buys in. Next turn it is \V{hivejack}ed while still at 2 \Vp{node}. As the only holder you collect \V{root} and
\V{sudo} (\Y{3,000}), then \V{cashout} your 4 \V{coin} for \Y{800}. You spent \Y{600} and got \Y{3,800} back.
\end{example}

\subsection{Middle: pick your moment}
\begin{itemize}
  \item \textbf{Look at the holdings.} \VC{coin} holdings are public. Before you \V{broker} a \V{hive} to a syndicate, or set off a
    \V{hivejack}, check who would be paid. Don't hand a rival a bonus.
  \item \textbf{One \V{coin} more than the leader takes \V{root}.} A tie for \V{root} splits the money. If you can see you are tied,
    a single extra \V{coin} takes the whole \V{root} bonus.
  \item \textbf{Grow it before it's \V{absorbed}.} Payouts and \V{cashout} prices use the \V{absorbed} \V{hive}'s size before the
    \V{hivejack}. If you hold its \V{coin}, push it up to the next step of the price chart first. Check the chart: a low-tier
    \V{hive} pays the same at 6 \Vp{node} as at 10, so growing it from 6 to 10 earns you nothing.
  \item \textbf{Don't grow a rival's \V{hive}.} Every \V{spread} raises the price and moves a \V{hive} toward a \V{hivejack} or
    \V{safe}. Make sure it is your \V{hive} you are helping.
  \item \textbf{Keep options in your \V{kit}.} Some \Vp{exploit} touch nothing and just pass the turn quietly. Others found, spread,
    or \V{hivejack}. Holding a mix lets you choose when things happen instead of being forced.
  \item \textbf{Count \Vp{exploit}.} There is exactly one \V{exploit} for each \V{node}. If a key \V{node} stays empty turn after turn,
    someone may be holding its \V{exploit} for the right moment.
\end{itemize}

\subsection{\VC{hivejack} choices}
\begin{itemize}
  \item \textbf{\VC{swap} when the \V{absorbed} \V{coin} was cheap.} A \V{swap} is always 2 for 1, whatever the sizes. Two \Vp{coin}
    worth \Y{300} each (\Y{600} if you \V{cashout}) for one worth \Y{800} is a good trade, and it might make you \V{root} of the
    \V{absorber}.
  \item \textbf{\VC{hodl} if you can bring it back.} If your \V{kit} can found a new \V{hive}, keeping \V{flatlined} \V{coin} and then
    brokering that syndicate again gives you a head start on \V{root}. Selling now and buying back later costs turns.
  \item \textbf{Cash early, \V{swap} late.} In the first few \Vp{hivejack}, \V{money} buys you more \V{coin}. Later, one \V{hive} usually
    gets big enough to last the game. Get into it with a \V{swap} before the end, or make sure you hold \V{root} somewhere else that is big.
\end{itemize}

\subsection{Late: the \V{nullroute}}
\begin{itemize}
  \item \textbf{\VC{safe} \Vp{hive} survive.} At \Rule{safesize} \Vp{node} a \V{hive} can never be \V{absorbed}, so its \V{root} bonus is
    settled at the final \V{cashout}. A \V{hive} one \V{node} short of \V{safe} is the best target on the \V{grid}, for you or for a rival.
  \item \textbf{\VC{flatlined} \V{coin} is worth nothing at the end.} \VC{hodl}ing late in the game is a gamble that the syndicate
    comes back in time.
  \item \textbf{Only call the \V{nullroute} if you are ahead.} Add up your \Game{currency}, your \V{coin} at current prices, and the bonuses
    you would collect. If someone else comes out on top, keep playing.
  \item \textbf{Use the other \Vp{player}' turns.} The shot clock only runs on your decisions, so plan your next \V{spike} while others
    move. With move hints on, each empty \V{node} your \V{kit} can hit shows the number of the \V{exploit} that hits it.
\end{itemize}

{\small\color{spike} Some of these tips are adapted from Robert Crowter-Jones's \emph{Acquire} strategy notes (Elusive Meeple, 2018).\par}
